\documentclass[11pt]{ctexart}
\usepackage[a4paper,margin=1in]{geometry}
\usepackage{graphicx}
\usepackage{xcolor}
\usepackage{hyperref}
\definecolor{refgray}{gray}{0.45}
\title{\textbf{市场最新观点汇总}\\[0.4em]{\large 结构性分化与供给侧再定价成为今日市场主线，AI与能源驱动出口窄脉冲，人民币走强信号被低估。}}
\author{KC桌面 自动生成}
\date{260523}
\begin{document}
\maketitle
\section*{一页摘要}
\begin{itemize}
  \item 中国WFE进口下滑主因光刻机供给瓶颈，非光刻设备需求依然强劲，结构性分化明显。
  \item 中国出口增长高度集中于AI存储芯片和新能源板块，58\%增量来自三个领域，价格驱动而非数量驱动。
  \item 日本股市在3\%利率环境下仍可维持上行，金融股相对受益，AI半导体行情不会因利率上升终结。
  \item 人民币走强基于贸易顺差、企业结汇韧性及地缘风险不对称，美元走强时抗跌、走弱时弹性更高。
  \item 美国变压器缺口收窄速度慢于预期，2028年仍将依赖进口满足六成需求。
  \item 股份制银行利润增速分化加剧，资产质量出清节奏差异是核心变量。
  \item 铝的供给侧拐点被低估，印尼产能释放时间差是关键变量，LME库存已至临界水平。
  \item 储能与新能源车领域成本通胀正在重新定义赢家，龙头集中度有望提升。
  \item 美国页岩油生产商维持资本纪律，产量仅增0.6\%，油价上涨现金流流向股东回报。
  \item 中国创新药在ASCO 2026的信号是Keytruda的“最佳助推器”出现，而非替代者。
\end{itemize}
\section{宏观与利率}
\textbf{全球利率环境分化，日本利率上升由结构性通胀预期驱动而非财政恐慌，人民币走强逻辑被市场低估。}

\begin{itemize}
  \item 日本10年期国债收益率逼近3\%关口，但利率上升主要源于盈亏平衡通胀率走高和进口价格回升，而非补充预算传闻，市场对财政恐慌的解读可能过度。
  \item 日本股市在3\%利率环境下仍能维持上升趋势，日元快速升值风险消退、区域金融体系韧性增强以及外资和日本个人投资者仍有买入空间是三大支撑。
  \item 人民币对美元呈现不对称反应：美元走强时人民币敏感度仅28\%，美元走弱时敏感度跃升至62\%，这种抗跌性源于持续的贸易顺差和企业结汇韧性。
  \item 人民币中间价模型预测较前一日下降370个基点，主要贡献来自卢布和欧元，模型误差持续缩小表明政策意图正从被动跟随转向主动管理预期。
  \item 中国外部资产负债表正在经历市场尚未充分定价的再平衡，即使美元阶段性反弹，人民币贬值空间也极为有限。
\end{itemize}
\begin{figure}[htbp]
\centering
\includegraphics[width=0.88\linewidth]{figures/fig_046.jpg}
\caption{Figure 4 - Figure 4: Japan 10Y Breakeven Inflation Rate (BEI)}
\end{figure}
\begin{figure}[htbp]
\centering
\includegraphics[width=0.88\linewidth]{figures/fig_047.jpg}
\caption{Figure 5 - Figure 5: Japan Import Price Index}
\end{figure}
\begin{figure}[htbp]
\centering
\includegraphics[width=0.88\linewidth]{figures/fig_048.jpg}
\caption{Figure 6 - Figure 6: Economic Surprise Index}
\end{figure}
\begin{figure}[htbp]
\centering
\includegraphics[width=0.88\linewidth]{figures/fig_049.jpg}
\caption{Figure 7 - Figure 7: Consumer Confidence Index}
\end{figure}
{\small\color{refgray}\textbf{References:} [R003] 日本股市真正的考验不是3\%利率，而是谁能在利率上升中重新定价; [R004] 人民币走强的逻辑，市场可能低估了供给侧的支撑力; [R005] 人民币中间价模型正在发出一个被市场低估的信号}

\section{AI/算力与半导体设备}
\textbf{中国WFE市场并非需求放缓，而是光刻机供给瓶颈下的结构性分化；存储器扩产成为新引擎，本土设备商替代速度被低估。}

\begin{itemize}
  \item 4月中国WFE进口27亿美元，同比下降3\%，但剔除光刻机后其他设备进口环比增长15\%、同比增长5\%，沉积设备进口同比增长36\%。
  \item 光刻机进口仅1.42亿美元，同比暴跌60\%，创近两年最低，但这是供给端约束而非需求萎缩，一旦DUV产能改善，被压抑的需求可能集中释放。
  \item 中国存储器IDM厂商长鑫存储和长江存储正加速扩产，计划2026年各新增一座晶圆厂，2027年再各增两座，AI驱动的存储器超级周期是核心驱动力。
  \item 2026至2028年中国WFE市场规模预测上调至580亿、670亿和770亿美元，累计上调超260亿美元，主要驱动力来自存储器资本开支而非逻辑芯片。
  \item 中国存储器厂商从洁净室建设到投产仅需1年，远快于全球同行的2-3年，运营效率差异被市场低估。
\end{itemize}
\begin{figure}[htbp]
\centering
\includegraphics[width=0.88\linewidth]{figures/fig_002.jpg}
\caption{EXHIBIT 2 - EXHIBIT 2: Month-over-month showed 12\% decline, mostly dragged down by lithography import. Total WFE imports to China MoM growth by segment (monthly) (USDmn)}
\end{figure}
\begin{figure}[htbp]
\centering
\includegraphics[width=0.88\linewidth]{figures/fig_004.jpg}
\caption{EXHIBIT 4 - EXHIBIT 4: The lithography import was \textbackslash{}\$142mn in Apr 2026, much weaker than last year Lithography imports to China (monthly) (USDmn)}
\end{figure}
\begin{figure}[htbp]
\centering
\includegraphics[width=0.88\linewidth]{figures/fig_007.jpg}
\caption{EXHIBIT 7 - EXHIBIT 7: Netherlands' share was flattish, Japan decreased by 3ppt, US+Singapore+Malaysia in aggregate grew slightly Total WFE imports to China by region EXHIBIT 8: Excluding lithography, there has been a gradual uptick in impor}
\end{figure}
\begin{figure}[htbp]
\centering
\includegraphics[width=0.88\linewidth]{figures/fig_064.jpg}
\caption{EXHIBIT 3 - EXHIBIT 3: We revised up 2026E/2027E/2028E by \textbackslash{}\$2.9/\textbackslash{}\$7.3/\textbackslash{}\$16.4 bn}
\end{figure}
\begin{figure}[htbp]
\centering
\includegraphics[width=0.88\linewidth]{figures/fig_065.jpg}
\caption{EXHIBIT 4 - EXHIBIT 4: We revised up our numbers mostly due to stronger CapEx in memory China WFE Market Share by End Market (USD bn)}
\end{figure}
{\small\color{refgray}\textbf{References:} [R001] 中国WFE进口的真正信号：不是需求放缓，而是供给瓶颈下的结构性分化; [R007] 中国半导体设备市场正在经历的不是周期，而是供给侧的重新定价}

\section{能源与大宗商品}
\textbf{铝的供给侧拐点被低估，美国页岩油资本纪律持续，变压器缺口收窄速度慢于预期，大宗商品定价权正从需求端向供给端转移。}

\begin{itemize}
  \item 铝的供需错配比市场认知更深：LME库存已降至30万吨临界水平，中国宏桥指出印尼产能释放需2.5年以上，市场对供应快速跟上的假设可能被证伪。
  \item 美国页岩油生产商在WTI站上88美元背景下仍维持资本纪律，2026年产量指引仅上调0.6\%，资本开支几乎不变，油价上涨带来的现金流流向股东回报。
  \item 美国电力变压器供需缺口将从当前的72\%收窄至2028年的60\%，即使考虑所有已宣布的本地扩产计划，仍有六成需求依赖进口。
  \item 中国变压器出口增速放缓至27\%但并非需求见顶，而是结构性换挡：非洲市场暴增681\%，北美增长66\%，资源向需求最紧迫的市场集中。
  \item 刚果金铜钴成本压力上升，硫酸占现金成本比例从10\%涨至20\%，柴油从低个位数涨至高个位数，资源国风险重新成为定价因子。
\end{itemize}
\begin{figure}[htbp]
\centering
\includegraphics[width=0.88\linewidth]{figures/fig_069.jpg}
\caption{Exhibit 1 - Exhibit 1: We continue to estimate that US power transformer demand/local supply gap would narrow from \$72\textbackslash{}\%\$ to \$60\textbackslash{}\%\$ in 2028E; while US PPI has stabilized at a high level, China export price is up \$+23\textbackslash{}\%\$ yoy (Feb-Apr avg, for 2}
\end{figure}
\begin{figure}[htbp]
\centering
\includegraphics[width=0.88\linewidth]{figures/fig_070.jpg}
\caption{Exhibit 2 - Exhibit 2: China total transformer export value grew \$27\textbackslash{}\%\$ in April (slowing vs \$56\textbackslash{}\%\$ in March) Exhibit 3: China transformer export value to the US grew +95\% yoy in April (vs 118\% yoy in March)}
\end{figure}
\begin{figure}[htbp]
\centering
\includegraphics[width=0.88\linewidth]{figures/fig_074.jpg}
\caption{Exhibit 8 - Exhibit 8: Key raw materials for transformer, GOES, saw prices inch down, while copper price had a sharp increase in 4Q25, and dipped in March 2026}
\end{figure}
\begin{figure}[htbp]
\centering
\includegraphics[width=0.88\linewidth]{figures/fig_081.jpg}
\caption{Exhibit 1 - Exhibit 1: Looking ahead to 2027, oil E\&Ps offer a median FCF yield of 12\% at strip (\textbackslash{}\textasciitilde{}\textbackslash{}\$75 WTI). US Majors and Canadian oil sands producers sit closer to 8\% 2027 FCF Yield Sensitivity}
\end{figure}
\begin{figure}[htbp]
\centering
\includegraphics[width=0.88\linewidth]{figures/fig_087.jpg}
\caption{Exhibit 8 - Exhibit 8: Our 2026 capex estimates remain roughly unchanged on average as most operators maintained their activity plans...}
\end{figure}
{\small\color{refgray}\textbf{References:} [R008] 美国变压器缺口收窄的速度，比市场想象的要慢; [R010] 铝的供给侧拐点被低估，这才是大宗商品当前最被忽视的结构性机会; [R012] 市场真正低估的不是油价，而是供给纪律的持续性}

\section{地产与消费}
\textbf{消费领域品牌分层加速固化，地产相关信贷需求疲软，银行利润增速分化揭示资产质量出清节奏差异。}

\begin{itemize}
  \item 618首日直播GMV显示品牌增长分层已演变为竞争格局固化：毛戈平实现10\%以上同比增长且折扣率持平，而巨子生物同比下滑约72\%，华熙生物下滑约68\%。
  \item 毛戈平的增长是“量价齐稳”结果，仅以5-7个SKU实现超越同行增长，反映高端国货赛道定价权正在兑现。
  \item 股份制银行利润增速从1\%到高单位数的裂口源于资产质量出清节奏差异：A银行不良贷款覆盖率超200\%，管理层对全年增长有信心；B银行仍处于零售资产质量恶化早期。
  \item 贷款需求呈现零售凉、对公热格局，4月信贷收缩部分源于季节性，央行未给硬指标，零售需求普遍疲软。
  \item 净息差普遍看到企稳信号，多家银行预计全年息差收窄幅度较2025年明显改善，部分银行目标同比稳定或微升。
\end{itemize}
\begin{figure}[htbp]
\centering
\includegraphics[width=0.88\linewidth]{figures/fig_115.jpg}
\caption{Exhibit 1 - Exhibit 1: Proya/MAOGEPING/Shiseido/Winona/Comfy/Herborist introduced more SKUs yoy, while QuadHA/Collgene/Biohyalux/Timage cut back in listings Number of SKUs in the top tier KOLs' beauty presale livestreaming Number of SKUs in}
\end{figure}
{\small\color{refgray}\textbf{References:} [R009] 股份制银行的分化时刻：利润增速从1\%到高单位数意味着什么; [R016] 618首日直播GMV的真正信号：不是大盘回暖，而是品牌分层正在加速固化}

\section{区域市场}
\textbf{日本股市在利率上升中结构性分化，韩国外卖市场稳定但脆弱，中国医疗设备市场部分触底但IVD定价风险才开始。}

\begin{itemize}
  \item 日本股市在3\%利率环境下金融股相对受益，AI半导体短期承压但不会终结，地产、电力和航空面临实质性压力。
  \item 韩国外卖市场Baemin与Coupang Eats的GTV份额稳定在65:35，但利润修复依赖平台费用和广告收入而非运营效率，可持续性存疑。
  \item 中国医疗设备市场呈现结构性触底：奥林巴斯消化内镜业务下滑幅度未恶化，泰尔茂神经血管业务通过渠道扩张实现量增覆盖价降。
  \item 体外诊断领域定价重构风险被严重低估，Sysmex中国区销售额暴跌38.5\%，管理层对统一检测定价政策充满不确定性。
  \item 中国医药行业正从成本优势驱动的外包服务商转向具备全球定价权的创新资产持有者，AI正在渗透研发全链条。
\end{itemize}
\begin{figure}[htbp]
\centering
\includegraphics[width=0.88\linewidth]{figures/fig_043.jpg}
\caption{Figure 3 - Figure 1: Japan and US 10Y Interest Rates}
\end{figure}
\begin{figure}[htbp]
\centering
\includegraphics[width=0.88\linewidth]{figures/fig_053.jpg}
\caption{Figure 12 - Figure 12: Capital Adequacy Ratios of Regional Banks and Shinkin Banks vs. JGB 10Y Yield}
\end{figure}
\begin{figure}[htbp]
\centering
\includegraphics[width=0.88\linewidth]{figures/fig_055.jpg}
\caption{Figure 16 - Figure 16: Semiconductor Stock Price Performance During Periods of Sharp Rate Increases}
\end{figure}
\begin{figure}[htbp]
\centering
\includegraphics[width=0.88\linewidth]{figures/fig_102.jpg}
\caption{EXHIBIT 2 - EXHIBIT 2: Korea's food delivery market is stabilizing, but regulatory risks around Coupang and any shift in Baemin's ownership could still reshape competition. Korea FD Market Evolution}
\end{figure}
\begin{figure}[htbp]
\centering
\includegraphics[width=0.88\linewidth]{figures/fig_103.jpg}
\caption{EXHIBIT 3 - EXHIBIT 3: Intense rider subsidy competition has compressed per-order profitability in Korea's food delivery market from 2023.}
\end{figure}
{\small\color{refgray}\textbf{References:} [R003] 日本股市真正的考验不是3\%利率，而是谁能在利率上升中重新定价; [R015] 韩国外卖市场的稳定只是表象，真正的变局来自监管与所有权重构; [R017] 中国医疗设备市场的真实分化：部分触底，但IVD的定价风险才刚刚开始; [R018] 中国医药行业正在经历的不是周期反弹，而是商业模式的结构性重写}

\section{风险偏好与政策信号}
\textbf{市场风险偏好结构性分化，政策信号从被动跟随转向主动管理，IFRS 18将重塑信息博弈规则。}

\begin{itemize}
  \item 人民币中间价模型显示政策意图正从被动跟随市场转向主动管理预期，模型预测的370个基点降幅本质是一篮子压力集中释放。
  \item IFRS 18将于2027年1月生效，强制要求从营业利润开始编制间接法现金流表，将终结管理层与投资者之间的信息不对称。
  \item 中国出口增长并非全面复苏，而是AI硬件周期和地缘政治驱动的能源安全焦虑共同制造的窄脉冲，对GDP和就业拉动有限。
  \item 投资者对中国市场态度审慎但挑剔，AI、科技和供应链龙头共识稳固，工业板块关注具备全球竞争力的出口导向型玩家。
  \item 储能需求增速放缓被误读，关键在于成本弹性而非需求总量，锂价上涨带来的利润率压力正推动市场聚焦龙头。
\end{itemize}
\begin{figure}[htbp]
\centering
\includegraphics[width=0.88\linewidth]{figures/fig_057.jpg}
\caption{EXHIBIT 1 - EXHIBIT 1: Accountants see what is visible (e.g. IFRS vs GAAP). IFRS vs GAAP International accounting}
\end{figure}
\begin{figure}[htbp]
\centering
\includegraphics[width=0.88\linewidth]{figures/fig_059.jpg}
\caption{EXHIBIT 4 - EXHIBIT 4: Falling in love with IFRS 18? It's the requirement to present the cash flow statement using the indirect method (i.e. starting with operating income rather than net income) that makes it attractive! Silhouette of a cou}
\end{figure}
\begin{figure}[htbp]
\centering
\includegraphics[width=0.88\linewidth]{figures/fig_032.jpg}
\caption{Figure 1 - Figure 1: New drivers of China exports \%pts contribution to total exports (US\textbackslash{}\$) \%oya growth}
\end{figure}
\begin{figure}[htbp]
\centering
\includegraphics[width=0.88\linewidth]{figures/fig_034.jpg}
\caption{Figure 3 - Figure 3: China export growth breakdown}
\end{figure}
{\small\color{refgray}\textbf{References:} [R002] 中国出口的真相：不是全面复苏，而是“AI+能源”的窄脉冲; [R005] 人民币中间价模型正在发出一个被市场低估的信号; [R006] 会计标准也能让人“爱上”？IFRS 18的真正影响，市场还远未定价; [R011] 储能与新能源车的真正分水岭：成本通胀正在重新定义赢家}

\section*{结语}
综合来看，今日市场主线围绕结构性分化与供给侧再定价展开：AI与能源驱动出口窄脉冲，人民币走强逻辑被低估，大宗商品定价权向供给端转移，消费品牌分层加速固化，区域市场内部差异显著。投资者需关注供给端约束的持续性以及政策信号从被动到主动的转变。
\vfill
{\small\color{refgray}Personal reading notes and learning share only. Not investment advice.}
\end{document}
